% SOURCE: docs/比赛材料/参赛材料准备清单.md + 项目介绍-面向小组.md §四–§八
\section{使用方法}

\subsection{系统启动}

本系统为纯客户端 Web 应用，使用时无需安装任何软件或配置运行环境。获取方式为自包含 HTML5 离线文件夹，
双击 \texttt{index.html} 即可在浏览器中运行，无需网络连接或本地服务器。

首屏加载约 3--5 秒（取决于设备性能），期间显示加载动画。加载完成后，系统以默认参数自动启动双摆仿真，
3D 场景淡入呈现。

\subsection{探索模式——感知混沌}

探索模式（快捷键 \textbf{1}）是系统的默认入口，提供最直观的双摆交互体验：

\begin{itemize}
  \item \textbf{3D 场景操作}：鼠标左键拖拽旋转视角，滚轮缩放，右键平移。右上角视角预设按钮支持三视角一键切换：
       “实验员视角”（正面平视）、“上帝视角”（斜上方俯视）、“混沌视角”（自动缓慢旋转）；
  \item \textbf{运动尾迹}：摆球后方拖曳的彩色尾迹以颜色编码速度——蓝色（静止）→ 绿色 → 黄色 → 红色（高速），
       粗细随摆球速度动态变化。尾迹可在“显示选项”面板中切换开关或调整持久度；
  \item \textbf{参数调节}：左侧参数面板提供 10 个物理参数的滑块与数值输入双模式。
       参数分为两组——连续参数（质量、摆长、重力加速度、阻尼）和初始条件参数（角度、角速度）。
       参数的具体范围与默认值见表\ref{tab:params}；
  \item \textbf{声音化引擎}：点击底部声音开关按钮（需用户手势激活 AudioContext），
       系统将摆的运动实时转化为声音。上摆角度映射音高（$\theta_1 \to 220 + |\dot{\theta}_2| \times 200\;\mathrm{Hz}$），
       两摆夹角映射和声（纯五度 = 3:2 频率比；混沌态混入白噪声）；
  \item \textbf{蝴蝶效应对比器}：点击“蝴蝶效应”按钮，界面切换为左右并排的双 Viewport 布局。
       左侧金色摆 A 与右侧紫色摆 B 初始角度仅差 $10^{-6\circ}$，由两个独立 Worker 同步积分。
       分离度超过 90° 时触发脉冲提示，直观展示初值敏感性的指数放大效应；
  \item \textbf{时间反演实验}：点击“时间反演”按钮，可选择“精确反演”模式（直接回放 RingBuffer 历史帧，
       零误差回溯）或“数值反演”模式（Worker 以 $-dt$ 重新积分 ODE，因浮点误差被混沌指数放大而逐渐偏离原路径）。
       底部 Canvas 实时绘制漂移距离曲线，配合三阶段叙事提示（重合期 $\to$ 分离期 $\to$ 发散期）。
\end{itemize}

\begin{table}[htbp]
  \centering
  \caption{物理参数范围与默认值}
  \label{tab:params}
  \begin{tabularx}{\textwidth}{l c c c l}
    \toprule
    \textbf{参数} & \textbf{符号} & \textbf{范围} & \textbf{默认值} & \textbf{单位} \\
    \midrule
    上摆质量 & $m_1$ & 0.1 -- 5.0 & 1.0 & kg \\
    下摆质量 & $m_2$ & 0.1 -- 5.0 & 1.0 & kg \\
    上摆杆长 & $L_1$ & 0.1 -- 3.0 & 1.0 & m \\
    下摆杆长 & $L_2$ & 0.1 -- 3.0 & 1.0 & m \\
    上摆初始角度 & $\theta_1(0)$ & $-\pi$ -- $\pi$ & 2.094 ($120^\circ$) & rad \\
    上摆初始角速度 & $\dot{\theta}_1(0)$ & $-10.0$ -- $10.0$ & 0.0 & rad/s \\
    下摆初始角度 & $\theta_2(0)$ & $-\pi$ -- $\pi$ & $-1.047$ ($-60^\circ$) & rad \\
    下摆初始角速度 & $\dot{\theta}_2(0)$ & $-10.0$ -- $10.0$ & 0.0 & rad/s \\
    重力加速度 & $g$ & 0.1 -- 20.0 & 9.81 & m/s$^2$ \\
    阻尼系数 & $b$ & 0.0 -- 2.0 & 0.0 & s$^{-1}$ \\
    \bottomrule
  \end{tabularx}
\end{table}

\subsection{分析模式——理解混沌}

分析模式（快捷键 \textbf{2}）提供定量分析工具，帮助学生从“观看”走向“分析”：

\begin{itemize}
  \item \textbf{Lyapunov 指数热力图}：在参数空间 $(L_2/L_1,\;\theta_{1,0})$ 上显示 $100 \times 100$ 网格的彩色热力图。
       红色区域（$\lambda_{\max} > 0$）对应混沌参数区，蓝色区域（$\lambda_{\max} \approx 0$）对应规则运动区。
       悬停任意单元格可查看 MLE 精确值；点击单元格可将对应参数组合一键注入仿真 Store，
       自动切换到探索模式观看该参数下的 3D 双摆运动——实现“从理论图谱到直观运动”的无缝跳转；
  \item \textbf{分岔图}：以 $\theta_{1,0}$ 为控制参数横轴，纵轴为稳态阶段 $\theta_2(t)$ 的局部极大值集合。
       周期运动对应有限个孤立点，混沌运动对应雾状散点带。支持 D3.js 缩放（框选放大）和竖直游标联动当前仿真参数；
  \item \textbf{庞加莱截面}：选取截面 $\theta_1 = 0$，穿越方向 $\dot{\theta}_1 > 0$（默认，可配置）。
       Worker 实时检测穿越事件并以线性插值定位精确穿越点，Canvas 2D 动态生长散点图，
       新点以 $\alpha$ 渐变加入，自动截断至 10,000 点以防止性能退化；
  \item \textbf{能量监控面板}：实时显示动能 $K$、势能 $V$、总机械能 $E = K + V$ 的时间曲线。
       能量相对漂移 $\eta(t) = |E(t)-E(0)|/|E(0)|$ 持续监控——$\eta > 0.5\%$ 时触发告警锁存（需用户手动清除），
       长时间仿真（$\geq$ 5 分钟）自动触发能量投影校正。
\end{itemize}

\subsection{实验模式——验证混沌}

实验模式（快捷键 \textbf{3}）提供物理正确性的自动化验证与高级分析工具：

\begin{itemize}
  \item \textbf{物理验证套件}：三项标准验证一键运行，每项拥有独立四项状态机
       （idle $\to$ running $\to$ passed $|$ failed）：
       \begin{enumerate}[label=(\alph*), leftmargin=*]
         \item \textbf{小角度线性化}：设置 $\theta_1 \approx 3^\circ,\;\theta_2 \approx 2^\circ$，
              对比 RKF45 数值解与线性化解析解的偏差，断言 $< 2\%$；
         \item \textbf{单摆退化}：设置 $m_2 \to 0$，测量摆动周期，断言与 $T_0 = 2\pi\sqrt{L_1/g}$ 偏差 $< 1\%$；
         \item \textbf{能量守恒}：关闭阻尼，运行 1000 s 仿真，断言能量相对漂移 $< 0.5\%$。
       \end{enumerate}
       全部通过显示绿色徽章，不通过自动诊断可能原因（步长过大 / 积分器精度不足 / 阻尼未关闭）；
  \item \textbf{受力拆解视图}：在 3D 场景中叠加矢量箭头，实时显示重力、上摆拉力、下摆拉力、切向惯性力、法向惯性力的方向与大小。
       支持笛卡尔坐标、极坐标、自然坐标三种坐标系切换；
  \item \textbf{Python 可编程沙箱}：点击“Lab 模式”，加载 Pyodide（浏览器内 Python 运行时）与 CodeMirror 6 编辑器。
       系统提供三个预设模板（弹簧摆 / 受迫摆 / 磁力摆）作为起点。用户修改运动方程后点击“运行”，
       Pyodide 调用 SciPy \texttt{solve\_ivp} 求解，结果通过 JsProxy 转换为 Float64Array 注入 3D 场景。
       代码执行超时（5 s）自动中断，Python 异常解析行号后映射为中文错误提示。
\end{itemize}

\subsection{故事模式——展示混沌}

故事模式（快捷键 \textbf{4}）为自动演示模式，适用于课堂展示或答辩场景。系统按预编排的 7 个阶段自动推进
（总时长约 3 分 30 秒），每阶段包含：参数渐变过渡（lerp）、相机程序化运动、电影式字幕淡入淡出、关键按钮脉冲高亮。
用户点击任意位置可暂停自动流程并恢复手动控制。

\subsection{辅助功能}

\begin{itemize}
  \item \textbf{状态快照}：点击快照按钮保存当前完整仿真状态（含参数、状态向量、尾迹数据和 64px 缩略图）至 IndexedDB，
       最多保留 50 个快照。快照卡片排列展示，点击恢复对应状态。支持双快照参数差异对比；
  \item \textbf{历史回放}：底部时间轴 Slider 基于 RingBuffer（100 s 历史），拖拽回溯任意历史时刻。
       支持分叉执行——在历史时刻 Fork 新仿真并以半透明轨迹并排对比；
  \item \textbf{实验报告导出}：一键生成 A4 PDF 报告，包含参数表、实验截图、能量曲线、混沌判定结论和误差分析表；
  \item \textbf{数据导出}：支持 CSV（轨迹时间序列）、JSON（完整状态快照）、PNG（Canvas 截图）格式导出，
       便于导入 Origin、MATLAB 等工具进行后续分析。
\end{itemize}

\endinput
